\subsection{COMPACT LIE GROUPS AND ROOT SYSTEMS}

With the terminology introduced in the preceding number, an important part of the results of numbers 4 and 5 can be summarized in the following theorem:

THEOREM 2. a) (X(T), X(T/(T ∩ D(G)), R(G, T))) is a reduced root diagram; its Weyl group consists of the X(w) for w ∈ W; the group X(C(G)) is isomorphic to the quotient of X(T) by the subgroup generated by R(G, T).

\begin{enumerate}
\def\labelenumi{\alph{enumi})}
\setcounter{enumi}{1}
\item
  \((\Gamma(\mathrm{T}),\Gamma(\mathrm{C}(\mathrm{G})_{0}),\mathrm{R}^{\vee}(\mathrm{G},\mathrm{T}))\) is a reduced root diagram; its Weyl group consists of the \(\Gamma(w)\) , for \(w\in W\) ; the group \(\pi_{1}(\mathrm{G})\) is isomorphic to the quotient of \(\Gamma(\mathrm{T})\) by the subgroup generated by \(\mathrm{R}^{\vee}(\mathrm{G},\mathrm{T})\) .
\item
  If each of the Z-modules X(T) and \(\Gamma(T)\) is identified with the dual of the other (no. 2, Prop. 3), each of the preceding root diagrams is identified with the inverse of the other.
\end{enumerate}

Denote by \(\mathrm{D}^{*}(\mathrm{G},\mathrm{T})\) the diagram \((\mathrm{X}(\mathrm{T}),\mathrm{X}(\mathrm{T}/(\mathrm{T}\cap\mathrm{D}(\mathrm{G})),\mathrm{R}(\mathrm{G},\mathrm{T})))\) and by \(\mathrm{D}_{*}(\mathrm{G},\mathrm{T})\) the diagram \((\varGamma(\mathrm{T}),\varGamma(\mathrm{C}(\mathrm{G})_{0}),\mathrm{R}^{\vee}(\mathrm{G},\mathrm{T}))\) ; these are called the contravariant diagram and the covariant diagram of G (relative to T), respectively.

Examples. 1) If G is semi-simple of rank 1, then \(\mathrm{D}^{*}(\mathrm{G},\mathrm{T})\) and \(\mathrm{D}_{*}(\mathrm{G},\mathrm{T})\) are necessarily isomorphic to one of the two diagrams \(\Delta_{2}=(\mathbf{Z},0,\{2,-2\})\) , \(\Delta_{1}=(\mathbf{Z},0,\{1,-1\})\) . If G is isomorphic to \(\mathbf{SU}(2,\mathbf{C})\) , \(\mathrm{D}_{*}(\mathrm{G},\mathrm{T})\) is isomorphic to \(\Delta_{1}\) (since G is simply-connected) so \(\mathrm{D}^{*}(\mathrm{G},\mathrm{T})\) is isomorphic to \(\Delta_{2}\) . If G is isomorphic to \(\mathbf{SO}(3,\mathbf{R})\) , \(\mathrm{D}^{*}(\mathrm{G},\mathrm{T})\) is isomorphic to \(\Delta_{1}\) (since \(\mathrm{C}(\mathrm{G})=\{1\})\) , so \(\mathrm{D}_{*}(\mathrm{G},\mathrm{T})\) is isomorphic to \(\Delta_{2}\) .

\begin{enumerate}
\def\labelenumi{\arabic{enumi})}
\setcounter{enumi}{1}
\item
  If G and \(G'\) are two connected compact Lie groups, with maximal tori T and \(T'\) , respectively, and if \(\mathrm{D}^{*}(\mathrm{G},\mathrm{T})=(\mathrm{M},\mathrm{M}_{0},\mathrm{R})\) and \(\mathrm{D}^{*}(\mathrm{G}',\mathrm{T}')=(\mathrm{M}',\mathrm{M}_{0}',\mathrm{R}')\) , then \(\mathrm{D}^{*}(\mathrm{G}\times\mathrm{G}',\mathrm{T}\times\mathrm{T}')\) can be identified with \((\mathrm{M}\oplus\mathrm{M}',\mathrm{M}_{0}\oplus\mathrm{M}_{0}',\mathrm{R}\cup\mathrm{R}')\) . Similarly for the covariant diagrams.
\item
  Let N be a closed subgroup of T, central in G, and let \((M, M_{0}, R)\) be the contravariant diagram of G relative to T. Then the contravariant diagram of G/N relative to T/N can be identified with \((M', M_{0}', R)\) , where \(M'\) is the subgroup of M consisting of the \(\lambda\) such that \(\lambda(N) = \{1\}\) and \(M_{0}' = M' \cap M_{0}\) .
\item
  Similarly, let N be a finite abelian group, and \(\varphi : \pi_{1}(G) \to N\) a surjective homomorphism. Let \(G'\) be the covering of G associated to this homomorphism; this is a connected compact Lie group, of which N is a central subgroup (General Topology, Chap. XI, in preparation), and G can be naturally identified with \(G'/N\) . Let \(T'\) be the maximal torus of \(G'\) that is the inverse image of T. If \((P, P_{0}, S)\) is the covariant diagram of G relative to T, the covariant diagram of \(G'\) relative to \(T'\) can be identified with \((P', P_{0}', S)\) , where \(P'\) is the kernel of the composite homomorphism \(\varphi \circ f(G, T) : P \to N\) (cf.~no. 6, Prop. 11), and \(P_{0}' = P_{0} \cap P'\) .
\end{enumerate}

Remarks. 1) Let \(\mathfrak{c}\) be the centre of \(\mathfrak{g}_{\mathbf{C}}\) ; then \(\mathfrak{c} = \mathrm{L}(\mathrm{C}(\mathrm{G}))_{(\mathbf{C})}\) . We have the following relations between the diagrams of G relative to T and the direct and inverse root systems of the split reductive algebra \((\mathfrak{g}_{\mathbf{C}}, t_{\mathbf{C}})\) :

\begin{enumerate}
\def\labelenumi{\alph{enumi})}
\tightlist
\item
  The canonical isomorphism from \(\mathbf{C} \otimes \Gamma(\mathrm{T})\) to \(\mathfrak{t}_{\mathbf{C}}\) induces a bijection from \(\mathbf{C} \otimes \Gamma(\mathrm{C}(\mathrm{G})_0)\) to \(\mathfrak{c}\) and a bijection from \(1 \otimes \mathrm{R}^{\vee}(\mathrm{G}, \mathrm{T})\) to \(2\pi i.\mathrm{R}^{\vee}(\mathfrak{g}_{\mathbf{C}}, \mathfrak{t}_{\mathbf{C}})\) .
\end{enumerate}

from \(\mathbf{C}\otimes\varGamma(\mathrm{C}(\mathrm{G})_{0})\) to c and a bijection from \(1\otimes\mathrm{R}^{\vee}(\mathrm{G},\mathrm{T})\) to \(2\pi i.\mathrm{R}^{\vee}(\mathfrak{g}_{\mathbf{C}},\mathfrak{t}_{\mathbf{C}})\) . b) The canonical isomorphism from \(\mathbf{C}\otimes\mathrm{X}(\mathrm{T})\) to the dual \(t_{C}^{*}\) of \(t_{C}\) induces a bijection from \(\mathbf{C}\otimes\mathrm{X}(\mathrm{T}/(\mathrm{T}\cap\mathrm{D}(\mathrm{G})))\) to the orthogonal complement of \(t_{C}\cap\mathcal{D}(\mathfrak{g})_{C}\) , and a bijection from \(1\otimes\mathrm{R}(\mathrm{G},\mathrm{T})\) to \(\mathrm{R}(\mathfrak{g}_{\mathbf{C}},\mathfrak{t}_{\mathbf{C}})\) .

\begin{enumerate}
\def\labelenumi{\arabic{enumi})}
\setcounter{enumi}{1}
\tightlist
\item
  Assume that the group G is semi-simple; denote by R (resp. \(R^{\vee}\) ) the root system R(G, T) (resp. \(R^{\vee}(G, T)\) ), so that we have inclusions
\end{enumerate}

\[
\mathrm{Q} (\mathrm{R}) \subset \mathrm{X} (\mathrm{T}) \subset \mathrm{P} (\mathrm{R}) \quad \mathrm{Q} (\mathrm{R} ^ {\vee}) \subset \Gamma (\mathrm{T}) \subset \mathrm{P} (\mathrm{R} ^ {\vee}).
\]

The finite abelian groups \(\mathrm{P}(\mathrm{R})/\mathrm{Q}(\mathrm{R})\) and \(\mathrm{P}(\mathrm{R}^{\vee})/\mathrm{Q}(\mathrm{R}^{\vee})\) are in duality (Chap. VI, §1, no. 9); if \(M^{*}\) denotes the dual group of the finite abelian group M, we deduce from the preceding canonical isomorphisms

\[
\begin{array}{l l} \Gamma (\mathrm{T}) / \mathrm{Q} (\mathrm{R} ^ {\vee}) \to \pi_ {1} (\mathrm{G}) & \mathrm{P} (\mathrm{R} ^ {\vee}) / \Gamma (\mathrm{T}) \to \mathrm{C} (\mathrm{G}) \\ \mathrm{P} (\mathrm{R}) / \mathrm{X} (\mathrm{T}) \to (\pi_ {1} (\mathrm{G})) ^ {\wedge} & \mathrm{X} (\mathrm{T}) / \mathrm{Q} (\mathrm{R}) \to (\mathrm{C} (\mathrm{G})) ^ {\wedge}. \end{array}
\]

In particular, the product of the orders of \(\pi_{1}(\mathrm{G})\) and \(\mathrm{C}(\mathrm{G})\) is equal to the connection index f of \(\mathrm{R}(\mathrm{G},\mathrm{T})\) (loc. cit.).

Now let \(G'\) be another connected compact Lie group, \(T'\) a maximal torus of \(G'\) . Let \(f: G \to G'\) be an isomorphism of Lie groups such that \(f(T) = T'\) ; denote by \(f_{T}\) the isomorphism from T to \(T'\) that it defines. Then \(\mathrm{X}(f_{\mathrm{T}})\) is an isomorphism from \(\mathrm{D}^{*}(\mathrm{G}', \mathrm{T}')\) to \(\mathrm{D}^{*}(\mathrm{G}, \mathrm{T})\) , denoted by \(\mathrm{D}^{*}(f)\) , and \(\Gamma(f_{\mathrm{T}})\) is an isomorphism from \(\mathrm{D}_{*}(\mathrm{G}, \mathrm{T})\) to \(\mathrm{D}_{*}(\mathrm{G}', \mathrm{T}')\) , denoted by \(\mathrm{D}_{*}(f)\) . If \(t \in T\) , and if we put \(g = f \circ \operatorname{Int} t = (\operatorname{Int} f(t)) \circ f\) , then \(\mathrm{D}^{*}(g) = \mathrm{D}^{*}(f)\) , \(\mathrm{D}_{*}(g) = \mathrm{D}_{*}(f)\) .

PROPOSITION 15. Let \(\varphi\) be an isomorphism from \(\mathrm{D}^*(\mathrm{G}',\mathrm{T}')\) to \(\mathrm{D}^*(\mathrm{G},\mathrm{T})\) (resp. from \(\mathrm{D}_*(\mathrm{G},\mathrm{T})\) to \(\mathrm{D}_*(\mathrm{G}',\mathrm{T}')\) ). There exists an isomorphism \(f:\mathrm{G}\to \mathrm{G}'\) such that \(f(\mathrm{T}) = \mathrm{T}'\) and \(\varphi = \mathrm{D}^*(f)\) (resp. \(\varphi = \mathrm{D}_*(f)\) ); if \(f_{1}\) and \(f_{2}\) are two such isomorphisms, there exists an element \(t\) of \(\mathrm{T}\) such that \(f_{2} = f_{1}\circ \operatorname {Int}t\) .

The second assertion follows immediately from Prop. 9 (no. 4); we prove the first for the covariant diagrams, for example. Denote by \(\mathfrak{g}'\) (resp. \(\mathfrak{t}'\) ) the Lie algebra of \(G'\) (resp. \(T'\) ), and by \(\mathfrak{g}_{\mathbf{C}}'\) (resp. \(\mathfrak{t}_{\mathbf{C}}'\) ) its complexified Lie algebra. By Chap. VIII, §4, no. 4, Th. 2 (i), there exists an isomorphism \(\psi : \mathfrak{g}_{\mathbf{C}} \to \mathfrak{g}_{\mathbf{C}}'\) that maps \(\mathfrak{t}_{\mathbf{C}}\) to \(\mathfrak{t}_{\mathbf{C}}'\) and induces on \(\Gamma(T) \subset \mathfrak{t}_{\mathbf{C}}\) the given isomorphism \(\varphi : \Gamma(T) \to \Gamma(T')\) . Then \(\mathfrak{g}\) and \(\psi^{-1}(\mathfrak{g}')\) are two compact forms of \(\mathfrak{g}_{\mathbf{C}}\) that have the same intersection \(\mathfrak{t}\) with \(\mathfrak{t}_{\mathbf{C}}\) ; by §3, no. 2, Prop. 3, there exists an inner automorphism \(\theta\) of \(\mathfrak{g}_{\mathbf{C}}\) inducing the identity on \(\mathfrak{t}_{\mathbf{C}}\) and such that \(\theta(\mathfrak{g}) = \psi^{-1}(\mathfrak{g}')\) . By replacing \(\psi\) by \(\psi \circ \theta\) , we can assume that \(\psi\) maps \(\mathfrak{g}\) to \(\mathfrak{g}'\) . Further, by Prop. 4 of no. 2, there exists a unique morphism \(f_{\mathrm{T}} : T \to T'\) such that \(\Gamma(f_{\mathrm{T}}) = \varphi\) . Then the restriction of \(\psi\) to \(\mathfrak{t}\) is \(L(f_{\mathrm{T}})\) , and by §2, no. 6, Prop. 8, there exists a unique morphism \(f : G \to G'\) that induces \(f_{\mathrm{T}}\) on \(T\) and \(\psi\) on \(g_{C}\) . Applying the preceding to \(\varphi^{-1}\) and \(\psi^{-1}\) we obtain an inverse morphism to f, which is therefore an isomorphism. Then \(\mathrm{D}_{*}(f)=\Gamma(f_{\mathrm{T}})=\varphi\) , hence the proposition.

Note that, if T and \(T'\) are two maximal tori of G, the diagrams \(\mathrm{D}^{*}(\mathrm{G},\mathrm{T})\) and \(\mathrm{D}^{*}(\mathrm{G},\mathrm{T}')\) are isomorphic (if \(g \in G\) is such that \(gTg^{-1} = T'\) , then \(\operatorname{Int} g\) is an isomorphism from G to G that maps T to \(T'\) ). Denote by \(\mathrm{D}^{*}(\mathrm{G})\) the isomorphism class of \(\mathrm{D}^{*}(\mathrm{G},\mathrm{T})\) (cf.~Theory of Sets, Chap. II, §6, no. 2); this is a root diagram that depends only on G and is called the contravariant diagram of G. The covariant diagram \(\mathrm{D}_{*}(\mathrm{G})\) of G is defined similarly, and we obtain:

COROLLARY. Two connected compact Lie groups G and \(G'\) are isomorphic if and only if the diagrams \(\mathrm{D}^{*}(\mathrm{G})\) and \(\mathrm{D}^{*}(\mathrm{G}')\) (resp. \(\mathrm{D}_{*}(\mathrm{G})\) and \(\mathrm{D}_{*}(\mathrm{G}')\) ) are equal.

PROPOSITION 16. For every reduced root diagram D, there exists a connected compact Lie group G such that \(\mathrm{D}^{*}(\mathrm{G})\) (resp. \(\mathrm{D}_{*}(\mathrm{G})\) ) is isomorphic to D.

\begin{enumerate}
\def\labelenumi{\alph{enumi})}
\item
  By replacing D, if necessary, by its inverse diagram, we are reduced to constructing G such that \(\mathrm{D}^{*}(\mathrm{G})\) is isomorphic to D. Put \(\mathrm{D} = (\mathrm{M}, \mathrm{M}_{0}, \mathrm{R})\) ; then \(Q \otimes M\) is the direct sum of \(Q \otimes M_{0}\) and the vector subspace \(\mathrm{V}(\mathrm{R})\) generated by R. Moreover, since the inverse roots take integer values on M, the projection of M on \(\mathrm{V}(\mathrm{R})\) parallel to \(Q \otimes M_{0}\) is contained in the group of weights \(\mathrm{P}(\mathrm{R})\) of R, so that M is a subgroup of \(\mathrm{M}_{0} \oplus \mathrm{P}(\mathrm{R})\) of finite index. Denote by \(D'\) the diagram \((\mathrm{M}_{0} \oplus \mathrm{P}(\mathrm{R}), \mathrm{M}_{0}, \mathrm{R})\) .
\item
  Let \(\mathfrak{a}\) be a complex semi-simple Lie algebra whose canonical root system is isomorphic to \(R \subset \mathbf{C} \otimes V(R)\) (Chap. VIII, §4, no. 3), and let \(\mathfrak{g}_1\) be a compact real form of \(\mathfrak{a}\) ( \(\S 3\) , no. 2, Th. 1). Let \(G_1\) be a simply-connected real Lie group whose Lie algebra is isomorphic to \(\mathfrak{g}_1\) ; then \(G_1\) is compact ( \(\S 1\) , no. 4, Th. 1). Let \(T_1\) be a maximal torus of \(G_1\) . By Th. 1, the diagram \(D^*(G_1, T_1)\) is isomorphic to \((P(R), 0, R)\) .
\item
  Let \(T_{0}\) be a torus of dimension equal to the rank of \(M_{0}\) ; then \(\mathrm{D}^{*}(\mathrm{T}_{0},\mathrm{T}_{0})\) is isomorphic to \((\mathrm{M}_{0},\mathrm{M}_{0},\varnothing)\) , so \(\mathrm{D}^{*}(\mathrm{G}_{1}\times\mathrm{T}_{0},\mathrm{T}_{1}\times\mathrm{T}_{0})\) is isomorphic to \(D'\) (Example 2).
\item
  Finally, let N be the finite subgroup of \(T_{1} \times T_{0}\) orthogonal to M. Put \(\mathrm{G} = (\mathrm{G}_{1} \times \mathrm{T}_{0})/\mathrm{N}\) , \(\mathrm{T} = (\mathrm{T}_{1} \times \mathrm{T}_{0})/\mathrm{N}\) . Then G is a connected compact Lie group, T a maximal torus of G, and D(G, T) is isomorphic to D (Example 3).
\end{enumerate}

Scholium. The classification of connected compact Lie groups up to isomorphism is thus reduced to that of reduced root diagrams. The connected compact semi-simple Lie groups correspond to the reduced root diagrams \((M, M_{0}, R)\) such that \(M_{0} = 0\) ; giving such a diagram is equivalent to giving a reduced root system R in a vector space V over Q and a subgroup M of V such that \(Q(R) \subset M \subset P(R)\) .

Remark 3. Let \(\mathrm{T}'\) be another maximal torus of G, B (resp. \(\mathrm{B}'\) ) a basis of the root system \(\mathrm{R}(\mathrm{G},\mathrm{T})\) (resp. \(\mathrm{R}(\mathrm{G}',\mathrm{T}')\) ) (Chap. VI, §1, no. 5, Def. 2). There exist elements \(g\in \mathrm{G}\) such that Int \(g\) maps T onto \(\mathrm{T}'\) and B onto \(\mathrm{B}'\) , and these elements form a unique coset modulo Int(T) (since T and \(\mathrm{T}'\) are conjugate, we can assume that \(\mathrm{T} = \mathrm{T}'\) , and it suffices to apply Chap. VI, §1, no. 5, Remark 4 and Prop. 9 of no. 4). It follows that the isomorphism from T to \(\mathrm{T}'\) induced by Int \(g\) is independent of the choice of \(g\) ; consequently the same is true of \(\mathrm{D}_{*}(\mathrm{Int}g)\) and \(\mathrm{D}^{*}(\mathrm{Int}g)\) . Paraphrasing Chap. VIII, §5, no. 3, Remark 2, mutatis mutandis, we can now define the canonical maximal torus of G, the canonical covariant and contravariant root diagrams of G, \ldots.

